% !TeX program = xelatex
% ================================================================
%  全国大学生数学建模竞赛（CUMCM）论文 LaTeX 模板
%  依据：《全国大学生数学建模竞赛论文格式规范》（2025 年修订稿）
%  参照：2023 年国赛 C 题论文排版
%
%  编译方式：XeLaTeX（务必用 XeLaTeX，中文支持最好）
%      命令行：xelatex main.tex   （需编译两次以生成交叉引用）
%      VS Code / TeXstudio 中请将编译引擎设为 XeLaTeX
%
%  【如何使用本模板】
%   1. 在本文中搜索 "【待填" 即可定位到每一处需要你填写/替换的位置。
%   2. 每个"【待填】"下方都写了：填什么、怎么填、注意事项。
%   3. 所有示例文字（如"本文针对……"）都只是占位，务必替换成你自己的内容。
% ================================================================

\documentclass[a4paper, zihao=-4, fontset=windows]{ctexart}

% ---------- 页面设置：A4，四边页边距 2.5cm（规范要求 ≥2.5cm）----------
\usepackage{geometry}
\geometry{a4paper, top=2.5cm, bottom=2.5cm, left=2.5cm, right=2.5cm}

% ---------- 行距：1.5 倍 ----------
\usepackage{setspace}
\onehalfspacing

% ---------- 页眉页脚：页码页脚居中，阿拉伯数字，从摘要页起第 1 页 ----------
\usepackage{fancyhdr}
\pagestyle{fancy}
\fancyhf{}                          % 清空所有页眉页脚
\fancyhead[L,R,C]{}                 % 页眉留空（不得出现身份信息）
\fancyfoot[C]{\thepage}             % 页码位于页脚中部
\renewcommand{\headrulewidth}{0pt}  % 去掉页眉横线

% ---------- 章节标题格式：黑体、居中 ----------
% ctex 默认编号：一级"一、"，二级"1.1"，三级"1.1.1"
\ctexset{
  section       = { format = {\heiti \zihao{3}  \centering}, beforeskip = 1.2em, afterskip = 0.8em },
  subsection    = { format = {\heiti \zihao{4}  \centering}, beforeskip = 1.0em, afterskip = 0.6em },
  subsubsection = { format = {\heiti \zihao{-4} \centering}, beforeskip = 0.8em, afterskip = 0.4em },
  abstractname  = {\heiti \zihao{-3}}
}

% ---------- 图表标题：表 N / 图 N，五号字，标签后空格 ----------
\usepackage{caption}
\captionsetup{labelsep=space, font=small, skip=6pt}

\usepackage{booktabs}      % 三线表（\toprule \midrule \bottomrule）
\usepackage{graphicx}      % 插图
\usepackage{float}         % [H] 强制固定位置
\usepackage{subcaption}    % 子图

% ---------- 数学 ----------
\usepackage{amsmath, amssymb}
\usepackage{amsthm}

% ---------- 算法伪代码 ----------
\usepackage{algorithm}
\usepackage{algpseudocode}
\floatname{algorithm}{算法}
\renewcommand{\algorithmicrequire}{\textbf{输入：}}
\renewcommand{\algorithmicensure}{\textbf{输出：}}

% ---------- 超链接（务必放在导言区最后）----------
\usepackage[colorlinks, linkcolor=black, citecolor=black, urlcolor=blue]{hyperref}

% ================================================================
%  正文
% ================================================================
\begin{document}

% ==================================================================
% 【待填 1】论文标题
% ------------------------------------------------------------------
%   填什么：把下面花括号里的「论文标题（请替换为您的论文题目）」
%           换成你的论文标题。
%   怎么填：标题应简洁凝练，一般写成「基于 XX 模型/算法的 XX 问题研究」，
%           能概括核心方法 + 研究对象即可，尽量不超过 30 字。
%   注意：绝对不能出现学校、姓名、队号等任何身份信息！
% ==================================================================
\begin{center}
  {\heiti \zihao{2} 论文标题（请替换为您的论文题目）}
\end{center}

\vspace{1.2em}

% ==================================================================
% 【待填 2】摘要正文
% ------------------------------------------------------------------
%   填什么：在 abstract 环境内，把下面四段示例文字换成你的真实摘要。
%   怎么填（国赛摘要通用写法，四段式）：
%       第一段：一句话点明问题背景 + 本文做了什么（总体思路）。
%       第二段：针对问题一，建立什么模型、用什么方法求解、得到什么结论。
%       第三段：针对问题二、问题三，依次概括模型与结论。
%       第四段：模型检验/灵敏度分析，点明模型的合理性与稳健性。
%   注意：
%       1. 摘要（含标题和关键词）必须控制在【一页以内】，不要超页！
%       2. 只写结论和数值结果，不要写推导过程。
%       3. 不能出现学校、姓名、队号等身份信息。
% ==================================================================
\begin{abstract}
本文针对……（第一段：概述问题背景与总体思路）。
本文首先……（第二段：针对问题一，建立……模型，采用……方法求解，得到……结论）。
其次……（第三段：针对问题二……）。再次……（针对问题三……）。
最后对模型进行……检验与灵敏度分析，验证了模型的合理性与稳健性。

\vspace{1em}

% ------------------------------------------------------------------
% 【待填 3】关键词
%   填什么：把「关键词一；关键词二；关键词三；关键词四」换成你的关键词。
%   怎么填：一般 3~5 个，用分号「；」分隔，词与词之间不要加标点结尾。
%   示例：补货定价；线性回归；整数规划；模拟退火算法
% ------------------------------------------------------------------
\noindent {\heiti 关键词：}关键词一；关键词二；关键词三；关键词四
\end{abstract}

\newpage

% ==================== 正文（从本页起连续编号）====================

% ==================================================================
% 【待填 4】问题重述
% ------------------------------------------------------------------
%   填什么：用自己的话复述赛题，先概括背景，再分条列出需要解决的问题。
%   怎么填：
%       - 开头一两段复述题目背景（不要照抄原题，要概括改写）。
%       - 下面的 enumerate 列表里，逐条列出题目要求你回答的几个问题。
%   注意：问题要完整、准确，不要遗漏题目中的子问题。
% ==================================================================
\section{问题重述}

\noindent
在现代生鲜蔬菜大型销售或仓储中，由于品种限制，不同产品虽有个体差异，但其对仓储和运输过程中的保鲜能力均提出了较高的要求……需要对生鲜商超的补货、仓储、销售等过程进行统筹规划，设计较为合理的管理和销售策略，从而降低仓储、运输和保存成本，实现总利润最大化。据此，本文需要解决以下问题：

\begin{enumerate}
  \item 问题一：……（复述问题一的具体要求）。
  \item 问题二：……（复述问题二的具体要求）。
  \item 问题三：……（复述问题三的具体要求）。
\end{enumerate}

% ==================================================================
% 【待填 5】问题分析
% ------------------------------------------------------------------
%   填什么：每个 subsection 下写一段，分析对应问题的本质、建模思路、求解方法。
%   怎么填：回答三个问题——「这道题到底在问什么」「打算用什么方法/模型」
%           「解题步骤是什么」。建议配一张总体思路流程图（见后面插图示例）。
%   注意：这里只讲思路，不要出现具体计算和结果。
% ==================================================================
\section{问题分析}

\subsection{对问题一的分析}

（分析问题一的本质、求解思路与方法。）对于问题一，……其本质是一个……问题，本文拟通过……方法建立……模型进行求解。

\subsection{对问题二的分析}

（分析问题二……。）

\subsection{对问题三的分析}

（分析问题三……。）

% ==================================================================
% 【待填 6】模型假设
% ------------------------------------------------------------------
%   填什么：列出建模前作的简化假设，一般 3~6 条。
%   怎么填：每条用 \item 写一句，用分号或句号结尾；假设要合理、可解释。
%   注意：假设不能与题目数据矛盾；假设是后面模型成立的前提，别乱写。
% ==================================================================
\section{模型假设}

为简化问题、抓住主要矛盾，本文作出如下假设：

\begin{enumerate}
  \item 假设题目所给数据真实、可靠；
  \item 假设不存在大型自然灾害（如干旱、洪涝）等不可抗力因素的影响；
  \item 假设市场金融货币没有大范围波动，价格与销量关系稳定；
  \item 假设……（补充其他必要假设）。
\end{enumerate}

% ==================================================================
% 【待填 7】符号说明（三线表）
% ------------------------------------------------------------------
%   填什么：把论文里用到的主要数学符号及其含义填进下面的三线表。
%   怎么填表格：
%       - 第 1 个花括号 {cl} 定义两列的列格式：c=居中、l=左对齐、r=右对齐。
%         本例 {cl} 表示第 1 列居中、第 2 列左对齐。
%       - 每一行里用 & 分隔单元格，行末用 \\ 换行。
%       - \toprule 表头顶线、\midrule 表头与内容分界线、\bottomrule 表底线。
%       - 公式用 $...$ 包起来（如 $x_{ij}$ 表示带下标的 x）。
%   注意：只列出正文里真正用到的符号，符号要与后文公式一致。
% ==================================================================
\section{符号说明}

本文用到的主要符号及其含义见表~\ref{tab:symbol}。

\begin{table}[htbp]
  \centering
  \caption{主要符号说明}
  \label{tab:symbol}
  \begin{tabular}{cl}
    \toprule
    符号 & 含义 \\
    \midrule
    $x_{ij}$ & 第 $i$ 种商品在第 $j$ 天的销售量 \\
    $p_i$   & 第 $i$ 种商品的定价 \\
    $c_i$   & 第 $i$ 种商品的单位成本 \\
    $N$     & 商品种类总数 \\
    $T$     & 决策周期的天数 \\
    \bottomrule
  \end{tabular}
\end{table}

% ==================================================================
% 【待填 8】模型的建立与求解（论文核心，篇幅最长）
% ------------------------------------------------------------------
%   填什么：按题目分成若干 subsection（问题一、问题二……），
%           每个问题下再细分 subsubsection，按"建模 → 求解 → 结果"展开。
%   怎么填：每建立一步都写明：为什么这样建模、公式含义、用了什么算法、
%           结果是什么。公式用 equation 环境，结果用三线表，趋势用图。
%   注意：
%       - 这是阅卷重点，推导要清楚，结果要与摘要一致。
%       - 表标题放在表【上方】，图标题放在图【下方】。
% ==================================================================
\section{模型的建立与求解}

% 下面三行是小节标题，把 \ldots 换成你问题一/二/三的实际主题。
\subsection{问题一：\ldots 的分布规律及相互关系}

\subsubsection{数据预处理}

（描述数据清洗、缺失值处理、异常值处理等。）本文首先对附件数据进行预处理，包括……。

\subsubsection{描述性统计分析}

（描述统计指标，给出表格，表标题在表上方。）由描述性统计得到表~\ref{tab:desc}：

% ------------------------------------------------------------------
% 【待填 9】示例：一张三线表（表标题在表上方）
%   怎么填：caption 里写表标题；{lcccc} 表示 5 列（1 列左对齐 + 4 列居中）；
%           把"品类 A/B/C"及数字换成你的真实数据，行数可增删。
% ------------------------------------------------------------------
\begin{table}[htbp]
  \centering
  \caption{各品类商品描述性统计一览表}
  \label{tab:desc}
  \begin{tabular}{lcccc}
    \toprule
    品类 & 均值 & 标准差 & 最大值 & 最小值 \\
    \midrule
    品类 A & 100 & 20 & 150 & 60 \\
    品类 B & 200 & 30 & 260 & 140 \\
    品类 C & 150 & 15 & 180 & 120 \\
    \bottomrule
  \end{tabular}
\end{table}

\subsubsection{相关性分析}

（建立相关模型，给出公式。）皮尔逊相关系数计算公式为
% ------------------------------------------------------------------
% 【待填 10】示例：公式
%   怎么填：用 equation 环境写公式；\label{eq:xxx} 起名，
%           正文里用 \eqref{eq:xxx}（或 \ref{eq:xxx}）引用它。
%           上下标、分式、求和等写法见下方公式示例。
% ------------------------------------------------------------------
\begin{equation}
  r = \frac{\sum_{i=1}^{n}(x_i-\bar{x})(y_i-\bar{y})}
           {\sqrt{\sum_{i=1}^{n}(x_i-\bar{x})^2}\sqrt{\sum_{i=1}^{n}(y_i-\bar{y})^2}},
  \label{eq:pearson}
\end{equation}
其中 $\bar{x}$、$\bar{y}$ 分别为变量 $x$、$y$ 的样本均值。

（给出图，图标题在图下方。）分析结果见图~\ref{fig:scatter}：

% ------------------------------------------------------------------
% 【待填 11】示例：插图（图标题在图下方）
%   怎么填：把 \fbox{...} 这一整行换成：
%       \includegraphics[width=0.7\textwidth]{图1.png}
%       - 图片文件（如 图1.png、scatter.pdf）要和 main.tex 放在同一文件夹。
%       - [width=0.7\textwidth] 表示宽度为版面的 70%，可改成具体数值如 8cm。
%   注意：建议图片用 PDF 或 PNG 格式；图标题写在 \caption 里，会自动编号。
% ------------------------------------------------------------------
\begin{figure}[htbp]
  \centering
  \fbox{\parbox[c][5cm][c]{0.7\textwidth}{\centering 此处插入图片，例如 \texttt{\textbackslash includegraphics[width=0.7\textbackslash textwidth]\{图1.png\}}}}
  \caption{各品类销量相关性散点图}
  \label{fig:scatter}
\end{figure}

\subsection{问题二：\ldots 与 \ldots 的关系}

（建立回归模型，给出公式。）由式~\eqref{eq:pearson} 的分析，建立线性回归模型
\begin{equation}
  y = kx + b,
  \label{eq:linear}
\end{equation}
其中 $k$ 为斜率，$b$ 为截距项。通过最小二乘法求解得到参数估计，结果见表~\ref{tab:reg}。

\begin{table}[htbp]
  \centering
  \caption{最小二乘回归参数统计表}
  \label{tab:reg}
  \begin{tabular}{lcc}
    \toprule
    品类 & 斜率 $k$ & 截距 $b$ \\
    \midrule
    品类 A & -2.137 & 11.453 \\
    品类 B & -0.028 & 1.247  \\
    品类 C &  1.247 & 5.766  \\
    \bottomrule
  \end{tabular}
\end{table}

\subsection{问题三：\ldots 的优化模型}

\subsubsection{目标函数的确立}

以总利润最大为目标，建立目标函数
\begin{equation}
  \max \; Z = \sum_{i=1}^{N}\sum_{j=1}^{T} \left( p_i - c_i \right) x_{ij},
  \label{eq:objective}
\end{equation}
其中各符号含义见表~\ref{tab:symbol}。

\subsubsection{约束条件的确立}

需求约束：各品类补货量应不低于每日销售量，即
\begin{equation}
  x_{ij} \ge 0, \qquad i=1,2,\dots,N,\; j=1,2,\dots,T.
\end{equation}

\subsubsection{模型求解}

采用模拟退火等智能算法求解上述优化模型，得到最优方案见表~\ref{tab:result}。

\begin{table}[htbp]
  \centering
  \caption{各品类定价与补货策略一览表}
  \label{tab:result}
  \begin{tabular}{lcc}
    \toprule
    品类 & 定价（元） & 补货量（kg） \\
    \midrule
    品类 A & 3.5 & 120 \\
    品类 B & 4.0 & 200 \\
    品类 C & 5.2 & 150 \\
    \bottomrule
  \end{tabular}
\end{table}

% ==================================================================
% 【待填 12】模型的评价与改进
%   填什么：优点、缺点、改进方向三部分，各写 2~4 条。
%   怎么填：优点写模型创新点；缺点客观承认简化假设带来的局限；
%           改进方向写"如果放宽某假设、换更优算法，结果会怎样"。
% ==================================================================
\section{模型的评价与改进}

\subsection{模型的优点}

\begin{enumerate}
  \item 模型……具有较强普适性；
  \item 采用……算法，求解效率高；
  \item ……。
\end{enumerate}

\subsection{模型的缺点}

\begin{enumerate}
  \item 模型在……方面作了简化假设；
  \item ……。
\end{enumerate}

\subsection{模型的改进方向}

（提出可改进方向，如引入更多影响因素、采用更精确算法等。）

% ==================================================================
% 【待填 13】参考文献（GB/T 7714 格式）
% ------------------------------------------------------------------
%   填什么：把你引用过的文献按顺序填进 \bibitem 里，删掉没用的示例。
%   怎么填：每条文献格式（文献类型标识见下方说明）：
%       专著  [M]：作者. 书名[M]. 版次. 出版地: 出版社, 年份.
%       期刊  [J]：作者. 题名[J]. 刊名, 年份, 卷(期): 起止页码.
%       学位论文[D]：作者. 题名[D]. 保存地: 保存单位, 年份.
%       标准  [S]：发布机构. 标准名称[S]. 年份.
%   注意：
%       1. 正文引用处要用 \cite{ref1} 标注（例如"见文献\cite{ref1}"）。
%       2. 只列正文真正引用了的文献，编号与 \cite 顺序一致。
% ==================================================================
\begin{thebibliography}{99}

\bibitem{ref1}
司守奎, 孙玺菁. 数学建模算法与应用[M]. 第 2 版. 北京: 国防工业出版社, 2015.

\bibitem{ref2}
姜启源, 谢金星, 叶俊. 数学模型[M]. 第 5 版. 北京: 高等教育出版社, 2018.

\bibitem{ref3}
全国大学生数学建模竞赛组委会. 全国大学生数学建模竞赛论文格式规范[S]. 2025.

\end{thebibliography}

\newpage

% ==================================================================
% 【待填 14】附录
%   填什么：
%       - 支撑材料文件列表：列出你随论文提交的支撑文件（程序、数据等）。
%       - 源程序代码：粘贴建模用到的全部可运行源代码。
%   注意：
%       1. 附录页数不限，但必须与正文装订在一起。
%       2. 若确实没有用程序，要在这里注明"本论文没有用到程序"。
% ==================================================================
\appendix

\section{支撑材料文件列表}

\begin{enumerate}
  \item \texttt{main.tex} —— 论文源文件；
  \item \texttt{code\_main.py} —— 建模主程序；
  \item \texttt{data\_processed.xlsx} —— 处理后的数据。
\end{enumerate}

\section{源程序代码}

% ------------------------------------------------------------------
% 【待填 15】源代码
%   怎么填：把下面 verbatim 环境里的示例代码，替换成你自己的完整可运行代码。
%   说明：verbatim 环境里的内容会原样输出（保留空格、换行），
%         适合贴 Python / MATLAB / Lingo 代码，中文注释也能正常显示。
%   注意：代码要完整、可运行，且运行结果要与论文一致，否则可能被取消评奖资格。
% ------------------------------------------------------------------
\begin{verbatim}
# -*- coding: utf-8 -*-
# 示例：皮尔逊相关系数计算
import numpy as np

def pearson(x, y):
    x = np.asarray(x, dtype=float)
    y = np.asarray(y, dtype=float)
    return np.corrcoef(x, y)[0, 1]

x = [1, 2, 3, 4, 5]
y = [2, 4, 6, 8, 10]
print(pearson(x, y))   # 输出 1.0
\end{verbatim}

\end{document}
